\section{Introduction}
\label{sec:introduction}

Modern connected vehicles are increasingly equipped with high resolution cameras, LiDAR, radar, telemetry modules, and driver or cabin monitoring systems that continuously observe both the external traffic scene and the internal vehicle context~\cite{rajashekara2024data,li2020smart}. Consequently, vehicles generate massive volumes of multimodal data, with recent surveys reporting data rates that can reach the terabyte per hour scale in advanced autonomous driving settings~\cite{rajashekara2024data,li2020smart}. Exploiting such data enables a broad class of learning tasks, including collision risk estimation, trajectory forecasting, cooperative perception, and personalized driver assistance, all of which require models that are both context aware and responsive to local driving conditions~\cite{yu2022personalized}.

Traditionally, such models have been trained by uploading vehicle data to cloud servers, where a global model is learned over the aggregated corpus. This paradigm is impractical in current vehicular environments as transferring massive amounts of data is prohibitively expensive both in terms of bandwidth and energy. Furthermore, it conflicts with strict latency constraints and automotive privacy regulations.
Federated learning (FL) partially addresses this tension by enabling collaborative training through model exchange rather than raw data sharing~\cite{mcmahan2017communication,lim2020survey}. Classical FL, however, remains fundamentally server centric. Its dependence on a central aggregator introduces a communication bottleneck and a single point of failure, while the learned objective is still biased toward a common model that may not adequately serve vehicles operating under distinct traffic regimes, road geometries, sensing conditions, and driver preferences.


To reduce dependence on remote cloud coordination, several vehicular learning systems have shifted aggregation to roadside units, base stations, or edge servers near the road network~\cite{wang2023mobilityaware,pervej2023resource,liao2025pfedvs}. Although this architecture lowers communication latency, it does not remove central orchestration. Vehicles must still maintain infrastructure connectivity, contend for shared uplink resources, and often contribute to regional or global models whose inductive bias remains insufficiently personalized. Mobility further complicates this design because short contact times, frequent handovers, and intermittent coverage can degrade both training efficiency and aggregation quality. More importantly, infrastructure centric coordination underuses a natural source of learning efficiency in vehicular environments, namely direct collaboration among vehicles that experience similar local contexts and can often exchange updates more cheaply than through a base station.

A more compelling substrate for collaboration is provided by hybrid 5G C-V2X networks, standardized by 3GPP~\cite{3gpp2021architecture}. In 3GPP compliant V2X systems, vehicles can communicate directly over the PC5 sidelink, operating in the 5.9\,GHz ITS band, to exchange messages directly with nearby vehicles at low latency and minimal infrastructure cost~\cite{molina2017ltev}. When longer-range communication is required, the Uu interface is used through the 5G cellular network at higher cost but broader reach~\cite{garcia2021tutorial}, thereby exposing two communication modes with markedly different latency, range, and resource profiles. This dual-interface capability creates a natural substrate for decentralized vehicle-to-vehicle (V2V) learning: nearby vehicles can collaborate efficiently and frequently share similar environment context, yet the most important collaborators may not always be the nearest ones. The PC5 sidelink can therefore support low latency exchange among adjacent vehicles, whereas the Uu path can selectively reach statistically useful but geographically distant collaborators through the cellular infrastructure. At the same time, local datasets are non-IID across vehicles, as each reflects a distinct driving environment. Effective collaboration must consequently be selective, personalized, and context aware.

Recent works in the collaborative learning literature~\cite{kharrat2025dpfl,liu2024dfedpgp,ye2023pfedgraph,chen2022graph,zhang2021fedfomo,ma2021pfedatt,zantedeschi2020fully} show that personalization and topology structure are key elements in achieving per client satisfactory models. Many approaches explored in the literature demonstrate that carefully selected collaborators can substantially improve vehicle specific performance under heterogeneous data and constrained communication. Yet the assumptions behind these approaches remain largely misaligned with vehicular networks. Most operate over comparatively stable client populations, infer collaboration in a logically centralized manner, or assume a single communication interface. Existing vehicular decentralized schemes exhibit the complementary limitation: they improve training efficiency under mobility and wireless constraints, but they do not jointly address personalization, hybrid link heterogeneity, or Byzantine robustness within a unified topology control problem~\cite{chen2025mobility,meng2025resource,yuan2023p2p}. In this setting, the collaboration topology is not a networking detail. It is a learning variable because the value of a collaborator depends jointly on vehicle similarity, environment context, link cost, and vehicle trustworthiness.

Vehicular neighborhoods are highly dynamic, as vehicles move rapidly across space and maintain only short-lived communication opportunities. Each vehicle must therefore learn under strict bandwidth, energy, and time constraints. Additionally, vehicular learning takes place in an open environment characterized by strong data heterogeneity and potentially unreliable participants. Personalization and Byzantine resilience are therefore essential requirements in V2X environments. Indiscriminate collaboration with many peers can waste scarce resources, increase computation overhead, and even degrade learning performance when selected neighbors are statistically irrelevant or malicious. These challenges motivate the central question of this paper: \emph{how should a vehicle in a hybrid 5G C-V2X network decide, at each stage of training, which neighbors to collaborate with and over which interface, so as to maximize personalized learning benefit under mobility, resource, and trust constraints without relying on centralized coordination?}


We answer this question by proposing \textsc{DANTE} (\textbf{D}ual-interface \textbf{A}ttention-guided \textbf{N}etwork-formation for \textbf{T}rustworthy vehicular \textbf{E}dge learning), a robust decentralized personalized learning framework built on dynamic network formation over a hybrid V2X collaboration graph. Unlike conventional vehicular FL, which treats communication as a mere transport layer, DANTE jointly models statistical relevance, radio reachability, and trust in topology control. Each vehicle maintains a personalized model and acts as a strategic agent that selects collaborators and assigns aggregation weights. The utility of each candidate link captures expected personalized loss reduction, communication and computation cost, and trust-aware robustness. Short-range exchanges are carried over the PC5 sidelink, while the Uu interface is used selectively to reach high-value collaborators beyond direct radio range. To adapt to rapidly changing neighborhoods, each vehicle learns its link formation policy through a local attention-guided PPO mechanism that encodes candidate neighbors, link quality, budget state, and past collaboration quality, and activates only links expected to improve personalization per unit cost. Training is fully decentralized: no central critic, shared parameters, or global state is required, and each agent learns solely from local observations and local reward. Robust aggregation and online trust estimation are integrated into the same loop, allowing vehicles to mitigate suspicious updates and revert to local training when collaboration becomes unreliable.

The main contributions of this paper are summarized as follows.
\begin{enumerate}
\item We formulate robust decentralized personalized learning over hybrid 5G C-V2X as a directed network formation game in which vehicles are strategic players, link formation decisions define the action space, and payoffs integrate personalized loss improvement, interface aware resource cost, and trust weighted collaboration quality.
\item Each vehicle independently learns a \textsc{Dante} policy for neighbor selection and resource allocation from local observations and local reward, without any central critic, shared parameters, or global state, consistent with the operational constraints of V2X networks.
\item We evaluate DANTE against several baselines within a high fidelity SUMO driven simulation testbed that reproduces real world urban road networks across six locations, spanning dense urban grids, high speed highway corridors, coastal boulevards, and campus road networks, under realistic time varying mobility and V2X link dynamics.
\end{enumerate}
